\begin{thebibliography}{99}
\footnotesize
\setlength{\itemsep}{1.5pt}\setlength{\parskip}{0pt}
\bibitem{jacovi2020} A.~Jacovi and Y.~Goldberg. Towards faithfully interpretable NLP systems: How should we define and evaluate faithfulness? \emph{Proceedings of the 58th Annual Meeting of the ACL}, pages 4198--4205, 2020.
\bibitem{molsanity} K.~Iyer. MolSanity: a reliability audit for feature attributions on molecular graph neural networks. \url{https://github.com/Kar488/molsanity}.
\bibitem{hoch1994} S.~J. Hoch, X.~Dr\`eze, and M.~E. Purk. EDLP, Hi-Lo, and margin arithmetic. \emph{Journal of Marketing}, 58(4):16--27, 1994.
\bibitem{ribeiro2016} M.~T. Ribeiro, S.~Singh, and C.~Guestrin. ``Why should I trust you?'' Explaining the predictions of any classifier. \emph{Proceedings of the 22nd ACM SIGKDD Conference}, pages 1135--1144, 2016.
\bibitem{lundberg2017} S.~M. Lundberg and S.-I. Lee. A unified approach to interpreting model predictions. \emph{Advances in Neural Information Processing Systems 30}, 2017.
\bibitem{lundberg2020} S.~M. Lundberg, G.~Erion, H.~Chen, A.~DeGrave, J.~M. Prutkin, B.~Nair, R.~Katz, J.~Himmelfarb, N.~Bansal, and S.-I. Lee. From local explanations to global understanding with explainable AI for trees. \emph{Nature Machine Intelligence}, 2(1):56--67, 2020.
\bibitem{adadi2018} A.~Adadi and M.~Berrada. Peeking inside the black-box: A survey on explainable artificial intelligence (XAI). \emph{IEEE Access}, 6:52138--52160, 2018.
\bibitem{adebayo2018} J.~Adebayo, J.~Gilmer, M.~Muelly, I.~Goodfellow, M.~Hardt, and B.~Kim. Sanity checks for saliency maps. \emph{Advances in Neural Information Processing Systems 31}, 2018.
\bibitem{kindermans2019} P.-J. Kindermans, S.~Hooker, J.~Adebayo, M.~Alber, K.~T. Sch\"utt, S.~D\"ahne, D.~Erhan, and B.~Kim. The (un)reliability of saliency methods. In \emph{Explainable AI: Interpreting, Explaining and Visualizing Deep Learning}, LNCS 11700, pages 267--280. Springer, 2019.
\bibitem{alvarezmelis2018} D.~Alvarez-Melis and T.~S. Jaakkola. On the robustness of interpretability methods. arXiv:1806.08049, 2018.
\bibitem{hooker2019} S.~Hooker, D.~Erhan, P.-J. Kindermans, and B.~Kim. A benchmark for interpretability methods in deep neural networks. \emph{Advances in Neural Information Processing Systems 32}, 2019.
\bibitem{rudin2019} C.~Rudin. Stop explaining black box machine learning models for high stakes decisions and use interpretable models instead. \emph{Nature Machine Intelligence}, 1(5):206--215, 2019.
\bibitem{rojat2021} T.~Rojat, R.~Puget, D.~Filliat, J.~Del~Ser, R.~Gelin, and N.~D\'iaz-Rodr\'iguez. Explainable artificial intelligence (XAI) on time series data: A survey. arXiv:2104.00950, 2021.
\bibitem{ismail2020} A.~A. Ismail, M.~Gunady, H.~Corrada~Bravo, and S.~Feizi. Benchmarking deep learning interpretability in time series predictions. \emph{Advances in Neural Information Processing Systems 33}, pages 6441--6452, 2020.
\bibitem{tonekaboni2020} S.~Tonekaboni, S.~Joshi, K.~Campbell, D.~Duvenaud, and A.~Goldenberg. What went wrong and when? Instance-wise feature importance for time-series black-box models. \emph{Advances in Neural Information Processing Systems 33}, pages 799--809, 2020.
\bibitem{montgomery1997} A.~L. Montgomery. Creating micro-marketing pricing strategies using supermarket scanner data. \emph{Marketing Science}, 16(4):315--337, 1997.
\bibitem{hoch1995} S.~J. Hoch, B.-D. Kim, A.~L. Montgomery, and P.~E. Rossi. Determinants of store-level price elasticity. \emph{Journal of Marketing Research}, 32(1):17--29, 1995.
\bibitem{abrams} E.~Abrams, G.~Gui, and A.~Horta\c{c}su. Finding exogenous variation in data. arXiv:1704.07787, 2017.
\bibitem{angrist2009} J.~D. Angrist and J.-S. Pischke. \emph{Mostly Harmless Econometrics: An Empiricist's Companion}. Princeton University Press, 2009.
\bibitem{bertrand2004} M.~Bertrand, E.~Duflo, and S.~Mullainathan. How much should we trust differences-in-differences estimates? \emph{Quarterly Journal of Economics}, 119(1):249--275, 2004.
\bibitem{cameron2015} A.~C. Cameron and D.~L. Miller. A practitioner's guide to cluster-robust inference. \emph{Journal of Human Resources}, 50(2):317--372, 2015.
\bibitem{dominicks} James M. Kilts Center for Marketing. Dominick's dataset. University of Chicago Booth School of Business.
\end{thebibliography}
